% ============================================================================
%  00-map.tex — bản đồ toàn cây english
%  Ngày tạo: 2026-09-06 | Sửa gần nhất: 2026-09-06 | Ôn gần nhất: chưa
%  Dependency: không. Mức độ: cốt lõi (đọc đầu tiên).
% ============================================================================
\documentclass[english.tex]{subfiles}
\begin{document}
\chapter*{Bản đồ cây --- quyển sách này là gì và đọc thế nào}
\addcontentsline{toc}{chapter}{Bản đồ cây}
\markboth{Bản đồ cây}{}

\begin{tomtat}
Đây là bản đồ của một quyển sách gồm 21 chương, viết cho một mục tiêu rất hẹp: \emph{hiểu đúng và sâu điều người viết hoặc người giảng thực sự muốn nói}, rồi diễn đạt lại ý của chính bạn cho chính xác. Nó không phải giáo trình luyện thi và không dạy ngữ pháp theo lối truyền thống --- ngữ pháp ở đây chỉ được dùng như công cụ phân tích câu. Năm phần: nền tảng về nghĩa (A), đọc sâu (B), viết rõ (C), sắc thái và ngữ dụng (D), và luyện hằng ngày (E). Nếu chỉ nhớ một điều: nghĩa của một từ không nằm trong từ điển mà nằm ở \emph{chỗ nó thường xuất hiện} --- và toàn bộ quyển sách này là các cách khai thác nhận xét đó.
\end{tomtat}

\section*{A. Quyển sách này giải quyết vấn đề gì}

Có một trạng thái rất quen với người đọc tài liệu chuyên ngành bằng tiếng Anh: bạn hiểu từng từ trong câu, tra được mọi từ lạ, nhưng đọc xong vẫn không nắm chắc tác giả đang khẳng định điều gì, khẳng định mạnh tới đâu, và điều gì họ cố tình \emph{không} nói. Bạn xem một khoá học và nghe rõ từng chữ, nhưng ra khỏi buổi học lại không tóm tắt được lập luận. Bạn viết một đoạn mô tả phương pháp, thấy nó ``không sai'' nhưng cũng không giống cách người bản ngữ trong ngành viết.

Ba tình trạng đó có chung một nguyên nhân, và nó không phải thiếu từ vựng theo nghĩa thông thường. Nguyên nhân là kiến thức về tiếng Anh được lưu dưới dạng \emph{cặp từ--nghĩa}: mỗi từ ứng với một nghĩa tiếng Việt. Cách lưu đó bỏ mất bốn thứ quyết định nghĩa thật: từ này thường đi với những từ nào, nó thuộc văn vực trang trọng hay đời thường, nó mang thái độ tích cực hay tiêu cực, và trong ngữ cảnh cụ thể này người viết dùng nó để \emph{làm gì}.

Quyển sách này được tổ chức quanh bốn thứ đó. Ngữ pháp chỉ chiếm một chương, và chỉ ở mức đủ để phân tích một câu dài thành các mệnh đề --- vì đó là mức thực sự cần cho việc đọc. Phần lớn trọng lượng dành cho ngữ nghĩa, ngữ dụng, cấu trúc lập luận, và kỹ thuật viết.

\section*{B. Năm phần và vì sao chia như vậy}

\begin{itemize}
\item \textbf{Phần A --- Nền tảng nghĩa (chương 1--4).} Trả lời câu hỏi ``nghĩa của một từ thực ra nằm ở đâu''. Gồm: quan hệ giữa từ, ngữ cảnh và ý định người nói; biết một từ là biết những gì; ngữ pháp tối thiểu để đọc đúng câu dài; và cấu tạo từ để đoán nghĩa từ mới.
\item \textbf{Phần B --- Đọc sâu (chương 5--10).} Trả lời câu hỏi ``làm sao lấy đúng ý từ một văn bản khó''. Ba tầng đọc; cấu trúc lập luận học thuật; phương pháp ba lượt cho bài báo khoa học; đọc giữa các dòng (hàm ý, hedging, thái độ); nghe giảng và xem khoá học; và cách xây từ vựng chuyên ngành.
\item \textbf{Phần C --- Viết rõ (chương 11--15).} Trả lời câu hỏi ``làm sao viết để người đọc hiểu đúng ý mình ngay lần đầu''. Câu rõ; mạch văn cũ-trước-mới-sau; đoạn, bài và tóm tắt; nói đúng mức bằng hedging; và quy trình sửa bài.
\item \textbf{Phần D --- Sắc thái và ngữ dụng (chương 16--18).} Phần khiến bạn nghe ra được điều người viết \emph{ngụ ý}: gần nghĩa không phải đồng nghĩa; kết hợp từ; và ẩn dụ định hình cách ngành bạn nói về vấn đề.
\item \textbf{Phần E --- Luyện hằng ngày (chương 19--21).} Bốn nhánh luyện của Nation; bộ công cụ (từ điển, kho ngữ liệu, thẻ ghi nhớ, mô hình ngôn ngữ); và lộ trình 12 tuần kèm ngân hàng tài nguyên.
\end{itemize}

\section*{C. Phụ thuộc và thứ tự đọc}

\begin{figure}[H]\centering
\begin{tikzpicture}[node distance=0.75cm and 1.25cm,
  p/.style={draw=steel,fill=bluebg,rounded corners=2.5pt,inner sep=6pt,align=center,
            font=\small,text width=3.15cm},
  hi/.style={draw=accent,fill=amberbg,rounded corners=2.5pt,inner sep=6pt,align=center,
            font=\small,text width=3.15cm},
  ar/.style={-{Latex[length=2.4mm]},draw=steel,thick},
  ard/.style={-{Latex[length=2.4mm]},draw=tiercol,thick,dashed},
  lb/.style={font=\scriptsize\itshape,color=reddark,align=center}]
\node[hi] (a) {\textbf{A.} Nền tảng nghĩa\\{\scriptsize ch.\ 1--4}};
\node[p,right=of a] (b) {\textbf{B.} Đọc sâu\\{\scriptsize ch.\ 5--10}};
\node[p,right=of b] (c) {\textbf{C.} Viết rõ\\{\scriptsize ch.\ 11--15}};
\node[p,below=1.35cm of b] (d) {\textbf{D.} Sắc thái và ngữ dụng\\{\scriptsize ch.\ 16--18}};
\node[hi,below=1.35cm of c] (e) {\textbf{E.} Luyện hằng ngày\\{\scriptsize ch.\ 19--21}};
\draw[ar] (a) -- node[lb,above=0pt,text width=2.2cm] {công cụ\\phân tích} (b);
\draw[ar] (b) -- node[lb,above=0pt,text width=2.2cm] {đọc nhiều\\thì viết được} (c);
\draw[ar] (c) -- (e);
\draw[ar] (a.south) |- node[lb,below left=2pt and 0.1cm,text width=2.4cm] {sắc thái là\\phần sâu của nghĩa} (d.west);
\draw[ard] (d.east) -- node[lb,right=1pt,text width=2.0cm] {đọc ra ngụ ý} (c.south);
\draw[ard] (d) -- (e);
\end{tikzpicture}
\caption{Quan hệ giữa năm phần. Hai phần tô hổ phách là phần \emph{khung}: A quyết định cách bạn hiểu mọi thứ sau đó, E quyết định bạn có duy trì được không. Đọc và viết nuôi nhau: bạn chỉ viết được cấu trúc mà bạn đã đọc thấy đủ nhiều, và bạn chỉ đọc nhanh ra cấu trúc khi đã từng tự viết nó.}
\label{fig:map-root-en}
\end{figure}

Ba lối đọc, tuỳ mục tiêu.

\emph{Mục tiêu chính là đọc tài liệu chuyên ngành:} chương 1--4, rồi phần B trọn vẹn, rồi chương 16--17. Phần C đọc sau, nhưng đừng bỏ --- viết lại một đoạn bằng lời của mình là bài kiểm tra đọc hiểu chặt nhất.

\emph{Mục tiêu chính là viết bài báo hoặc tài liệu kỹ thuật:} chương 1--2, rồi phần C trọn vẹn, rồi chương 6, 8, 14. Phần D đọc song song vì nó quyết định bạn chọn từ nào.

\emph{Muốn hiểu tiếng Anh sâu như một hệ thống:} đọc tuần tự A \(\rightarrow\) B \(\rightarrow\) C \(\rightarrow\) D \(\rightarrow\) E, mỗi tuần một chương, và làm bài tập của chương đó trong tuần.

\section*{D. Cốt lõi và tham khảo}

\textbf{Cốt lõi --- phải nắm chắc:} chương 1 (nghĩa nằm ở đâu), chương 2 (biết một từ là biết gì, và ngưỡng độ phủ 98\%), chương 6 (cấu trúc lập luận), chương 8 (hedging và hàm ý --- chương quan trọng nhất cho mục tiêu ``hiểu đúng ý người nói''), chương 11--12 (câu rõ và mạch văn), chương 17 (collocation), chương 19 (bốn nhánh luyện).

\textbf{Tham khảo --- biết là có, tra khi cần:} chi tiết hình thái học ở chương 4; phần lý thuyết thể loại ở chương 13; phần ẩn dụ ý niệm ở chương 18 đọc một lần cho biết cách nhìn, không cần thuộc.

\section*{E. Quy ước dùng trong cả quyển}

Mọi ví dụ tiếng Anh được in \emph{nghiêng}; phần dịch hoặc giải thích đi ngay sau bằng tiếng Việt. Khi so sánh cách viết, sách dùng khối \textbf{Trước và sau}: câu gốc, câu đã sửa, và một dòng nói \emph{vì sao} --- dòng này mới là phần đáng đọc.

Các khối lặp lại ở mọi chương: \textbf{Tóm tắt}, \textbf{Kiến thức cần có trước}, \textbf{Câu hỏi mở đầu} (được trả lời ở cuối chương), \textbf{Gốc rễ} (khái niệm này ra đời từ nhu cầu nào), \textbf{Liên hệ ngang} (đối chiếu với tiếng Việt và với các ngành khác), \textbf{Nói lại thật đơn giản}, \textbf{Active recall}, \textbf{Câu hỏi ôn tập có đáp án}, và \textbf{Bài tập} --- phần bài tập luôn có ít nhất một việc làm được ngay trong ngày.

Ba loại hộp màu dùng tiết chế: \textbf{Trực giác}, \textbf{Cạm bẫy} (lỗi mà người Việt học tiếng Anh hay mắc), \textbf{Cần nhớ}.
\end{document}
